\scanpage{63}\mspage{10}
% archive foliation: 62
une application canonique
\[
\tag{30}\Sigma_{\underline J}\xrightarrow{f_{\unsure{2}}^t}\Sigma_{\underline J}
\]
\margin{c'est l'application déduite par passage aux quotients de $2\operatorname{id}_{E_0}:E_0-E_0$}
qu'on doit pouvoir se définir en termes de $\underline J$, et de la seule section $t$ de $\Sigma_{\underline J}$ (ce qui en fait détermine $E_0$, à isom \inserted{(non unique)} près loc. pour la top. étale\ldots). On a i.e.
\margin{il faudrait pouvoir \unsure{donner une} description intrinsèque de $\unsure{f_2^t|f_i^t}$, sans faire intervenir des courbes elliptiques}
\deleted{$\unsure{I}=\underline J\amalg$}
\[
\tag{31}\underline I_t=\underline J\amalg t(S)
\]
\margin{\uline{NB} comme $S$-schéma, $\underline I_t$ ne dépend pas de $t$, mod isom. can.}
et \deleted{$\unsure{E_0}$} $\underline V$ peut se décrire comme $V_{\underline I}\simeq V_{\underline J}$, opérant sur $\underline I$ de la façon canonique, cette opération se prolongeant en une opération sur $\Sigma_{\underline J}$ (dépendant de la section $t$ de $\Sigma_{\underline J}^*$). Le \deleted{\unsure{applic}} morphisme $f_t$ est défini sans ambiguïté par les conditions suivantes (où on suppose pour simplifier que $\underline J=J_S$, $J$ \unsure{ens|un} \unsure{à|:} trois éléments):
\[
\tag{32}\left\{\begin{array}{l}
\forall i\in J,\ \text{si }\sigma_i^t\in\operatorname{Aut}_S(\Sigma_J)\text{ est l'autom. qui échange}\\
t\text{ et }i,\text{ et }j\text{ et }k\quad(\text{où }\{i,j,k\}=J),\text{ et}\\
\deleted{\text{alors }\unsure{\text{si }(\Sigma_J)^{\sigma_i}}}\quad\Delta_i^t=(\Sigma_J)^{\sigma_i^t},\text{ on a}\\
f(\Delta_i^t)=\text{section }\xi_i\text{ de }D_{\underline J}\simeq J_S
\end{array}\right.
\]
et cela doit impliquer automatiquement
\[
\tag{33}f_t(\unsure{I_S|I_t}=J_S\cup t(S))\subset\deleted{\text{section}}t(S)
\]
